\titledquestion{Problem Analysis}

Consider the following basic problem:

You're given an array $A$ consisting of n integers $A[1], A[2], ..., A[n]$. You'd like to output a two-dimensional $n$-by-$n$ array $B$ in which $B[i, j]$ (for $i < j$) contains the sum of array entries $A[i]$ through $A[j]$ -- that is, the sum $A[i] + A[i+1] + ... + A[j]$. (the value of array entry $B[i, j]$ is left unspecified whenever $i \ge j$.)

Here's the pseudocode of a simple algorithm to solve this problem. (Initially $B[i,j]=0$ for all valid $i,j$)

\begin{algorithmic}[1]
\For{$i = 1$ \textbf{to} $n$}
    \For{$j = i+1$ \textbf{to} $n$}
        \For{$k = i$ \textbf{to} $j$}
            \State $B[i,j]\gets B[i,j] + A[k]$
        \EndFor
    \EndFor
\EndFor
\end{algorithmic}

For the following questions,
\begin{itemize}
    \item Define $T(n)$ as the total running time of the algorithm.
    \item Suppose each operation like $B[i,j]\gets B[i,j] + A[k]$ costs a constant time $c$.
    \item Ignore the time that \lstinline|for| loop iterations take.
\end{itemize}
\begin{parts}
    \part[1] Give out a function $f$ satisfying that the running time of the algorithm is $\Theta(f(n))$. That is, $T(n)=\Theta(n^3)$. (Give your answer in the most simplified form.)
    \part[3] For this function chosen \( f \), show that the running time of the algorithm on an input of size \( n \) is asymptotically upper bounded by a constant multiple of \( f(n) \). That is, prove $T(n)=O(f(n))$.
    \begin{solution}
        $$
        \begin{aligned}
        T(n)
        &= c \sum_{i=1}^{n} \sum_{j=i+1}^{n} \sum_{k=i}^{j} \\
        &= c \sum_{i=1}^{n} \sum_{j=i+1}^{n} (j-i+1) \\
        &< c \sum_{i=1}^{n} \sum_{j=1}^{n} n \\
        &= c n^3 
        \end{aligned}
        $$
        
        So

        $$
        T(n) = O(n^3)
        $$
    \end{solution}
    \part[3] For the same function \( f \), show that the running time of the algorithm on an input of size \( n \) is also asymptotically lower bounded by a positive constant multiple of \( f(n) \). That is, prove $T(n)=\Omega(f(n))$.
    \begin{solution}
        $$
        T = \sum_{i=1}^{n-1} \sum_{j=i+1}^{n} (j - i + 1)
        $$

        $$
        \begin{aligned}
        \sum_{j=i+1}^{n} (j - i + 1)
        &= \sum_{m=1}^{n-i} (m + 1) \\
        &= \sum_{m=1}^{n-i} m + \sum_{m=1}^{n-i} 1 \\
        &= \frac{(n-i)(n-i+1)}{2} + (n-i) \\
        &= \frac{(n-i)(n-i+3)}{2}
        \end{aligned}
        $$

        $$
        \begin{aligned}
        T
        &= \sum_{i=1}^{n-1} \frac{(n - i)(n - i + 3)}{2} \\
        &= \frac12 \sum_{t=1}^{n-1}t^2 + \frac32 \sum_{t=1}^{n-1} t \\
        &= \frac{n(n-1)(n+4)}{6}
        \end{aligned}
        $$

        So

        $$
        T(n) = \Theta(n^3)
        $$

        As a result,

        $$
        T(n) = \Omega(n^3)
        $$

    \end{solution}
    \part[3] Although the algorithm you just analyzed is the most natural way to solve the problem---after all, it just iterates through the relevant entries of the array \( B \), filling in a value for each---it contains some highly unnecessary sources of inefficiency. Give a different algorithm to solve the problem, with an asymptotically better running time. In other words, you should design another algorithm with running time \( \Theta(g(n)) \), where \( \lim\limits_{n \to \infty} \frac{g(n)}{f(n)} = 0\).
    
    Write pseudocode and state its tight asymptotic running time.
    \begin{solution}
        \begin{lstlisting}[language=C++]
    for(int i = 1; i < n; i++) {
        int current_sum = A[i];
        for(int j = i + 1; j < n; j++) {
            current_sum = current_sum + A[j];
            B[i, j] = current_sum;
        }
    }
        \end{lstlisting}

        Use prefix sum to optimize cost $\Theta(n)$, and the rest part cost $\frac{n(n-1)}{2}=\Theta(n^2)$.

        As a result, the total running time is $\Theta(n^2)$.
    \end{solution}
\end{parts}